\chapter{Conclusion and Future Work}
\label{chap:conclusion}

\section{Summary}

This report set out to answer a concrete systems question: given a
heterogeneous data pool mixing structured columns and free-text reviews,
how should a query engine execute a natural-language question that
mixes a structured predicate with a semantic one, while controlling both
LLM invocation cost and wall-clock latency? We benchmarked two baseline
strategies -- row-wise predicate pushdown (SUQL~\cite{suql}) and an
adaptive, join-based batched strategy in the style of Trummer -- and
found them to occupy different points of the cost/quality trade-off
space, now validated over ten held-out, genre-diverse questions rather
than one (Chapter~\ref{chap:experiments}). We introduced a two-level
cascade, applicable on top of either baseline, in which a cheap
first-stage scorer resolves the majority of units of work immediately
and escalates only genuinely ambiguous units to an expensive verification
step. The cascaded batched design, Trummer~V1, dominated on every axis
we measured: best F1 ($0.757$), best precision ($0.810$), the highest
\emph{floor} across the ten questions ($0.500$, versus $0.267$ or lower
for every other method), and by far the lowest cost ($10.0$ calls,
$35.05$s per question) -- at a small ($\approx 7\%$ relative) recall cost
compared to the cascaded row-wise design. New in this revision, we also
proved -- rather than only observed -- \emph{why}: a calibrated cascade's
expected and (with high, exponentially-growing probability) realized
cost is lower than an all-expensive baseline's once the candidate pool
is large enough, provided the cheap stage's batching factor and
calibrated escalation rate satisfy an explicit margin condition
(Section~\ref{sec:theory-cost}). Plugging in the measured system
constants shows SUQL~V1 fails that margin condition on both calls and
time (explaining its wall-time regression analytically, not just
empirically), while Trummer~V1 satisfies it comfortably on both.

\section{Engineering Lessons}

Beyond the headline numbers, four lessons from this project generalize
beyond this specific workload:
\begin{enumerate}
  \item \textbf{Predicate pushdown is the dominant cost lever} --
        confirmed, not just asserted, by the fact that the one method in
        the Trummer family that adds it (Trummer~V1) is also the only
        one that is both fast \emph{and} accurate.
  \item \textbf{Batching the cheap stage is not an optimization detail
        -- it determines whether cascading can win on cost or time at
        all.} Section~\ref{sec:theory-cost}'s worked instantiation shows
        this precisely: SUQL~V1's unbatched ($B=1$) cheap stage cannot
        satisfy the theorem's margin condition under its current
        per-call latency, while Trummer~V1's batched ($B=8$) cheap stage
        satisfies it with room to spare.
  \item \textbf{Calibration conventions are a real source of bugs.} The
        $\alpha$-as-confidence-level vs.\ $\alpha$-as-tail-probability
        pitfall of Section~\ref{sec:quantile-pitfall} produced two
        individually defensible formulas, only one of which matched the
        calling code's own convention; only a regression test against a
        synthetic confusion matrix with a known closed-form answer
        caught it.
  \item \textbf{A single hand-verified query is not enough to trust a
        cost/quality comparison.} Moving from one query to the
        \texttt{10q} suite changed the story in a material way: it
        surfaced Trummer~V1's floor-robustness (invisible in a
        single-query comparison, since there is no floor to speak of
        with $n=1$) and it exposed a genuinely hard question (Q3,
        romance chemistry) on which two of the four methods nearly
        collapse -- exactly the kind of finding a broader benchmark
        suite is for.
\end{enumerate}

\section{Current Implementation Status}

At the time of writing, all four methods -- SUQL baseline/V1, Trummer
baseline/V1 -- are implemented, documented, and benchmarked end to end
in the public repository
(\url{https://github.com/AnnaRemi/ENSIMAG-AI-M2-final-project}, branch
\texttt{agent/consolidate-four-method-benchmarks}), with a shared runner
and evaluator (\texttt{benchmarks/shared/scripts/}) producing the
\texttt{aggregate.csv}, \texttt{comparison.csv}, and per-question plots
this report draws on. Four benchmark suites of increasing size
(\texttt{1q}, \texttt{3q}, \texttt{5q}, \texttt{10q}) are available and
runnable locally (against a local Ollama instance) or on the lab's
shared Aker cluster.

\section{Future Work}

\subsection{Batch SUQL~V1's cheap stage}
\label{sec:future-suql-v1}
Section~\ref{sec:theory-cost}'s analysis makes this the single most
consequential remaining change: SUQL~V1's cheap stage currently scores
one row per call ($B=1$), and its measured per-row latency already
exceeds the expensive stage's own per-call latency, which
Proposition~\ref{prop:crossover} shows makes a wall-time win
\emph{structurally impossible} regardless of calibration quality or
pool size. Batching the cheap probe (scoring several rows' log-odds in
one completion, the same lever that already makes Trummer~V1's cheap
stage effective) is predicted, by the same analysis, to be the change
most likely to convert SUQL~V1 from a quality-only cascade into one that
also wins on cost and time.

\subsection{Port Beta-posterior calibration to Trummer~V1}
As flagged in Section~\ref{sec:trummer-v1}, Trummer~V1's threshold is
currently learned by a point-estimate agreement search rather than the
Beta-posterior credible-bound procedure used by SUQL~V1
(Section~\ref{sec:beta-posterior}). The two are not equally rigorous:
only the Beta-posterior version currently carries an explicit,
finite-sample confidence-level guarantee. Unifying both cascades on the
same calibration procedure would let Trummer~V1 -- already the strongest
variant on every other axis -- also carry that guarantee.

\subsection{Amortize the calibration budget across repeated queries}
Both cascades currently re-run their calibration procedure fresh for
every question ($K_{\mathrm{calls}}=20$ row-wise calls for SUQL~V1,
$1$ batched call for Trummer~V1, per Table~\ref{tab:theory-constants}).
In a deployment that re-asks a stable set of semantic predicates
repeatedly (\eg a fixed catalogue of recommendation filters), this cost
is payable once and amortized over every future invocation of the same
predicate, pushing the crossover point $n^\ast$ of
Proposition~\ref{prop:crossover} even lower than the already-small value
computed for Trummer~V1 in Section~\ref{sec:worked-instantiation}.

\subsection{The KV-cache-aware execution ladder}
\label{sec:future-kv-cache}
Stretto's second contribution -- \emph{Expected Attention Press}, a
cache-aware execution ladder that reuses attention key/value state
across cascade stages -- remains deliberately deferred. It requires
direct access to the serving engine's attention state (\eg via
\texttt{vllm}), which an Ollama-served backend does not expose; adopting
it is future work contingent on a backend migration, and remains the
single largest piece of remaining architectural work identified by this
project.

\subsection{Extending the selectivity/difficulty grid}
\label{sec:future-ablation}
The \texttt{10q} suite fixes structural selectivity at $40\%$
(Chapter~\ref{chap:experiments}); the repository's smaller \texttt{1q},
\texttt{3q}, and \texttt{5q} suites, and a still-larger suite beyond
\texttt{10q}, are the natural next steps toward a full grid crossing
structural selectivity with semantic-predicate difficulty, which would
let the cost/quality trade-offs of this report be stated as trends
across that grid rather than at one fixed selectivity.

\subsection{Early termination for \texttt{LIMIT} queries}
As discussed in Chapter~\ref{chap:experiments}, wiring early stopping
into the cascade's main accept loop once $|\mathit{Accept}| \ge k$ for a
\texttt{LIMIT k} query remains a small, self-contained, and still
unimplemented optimization.
